\documentclass[11pt]{article}
\usepackage[paperwidth=5.6in,paperheight=3.6in,margin=0.35in]{geometry}
\usepackage{amsmath,amssymb}
\pagestyle{empty}
\setlength{\parindent}{0pt}
\setlength{\parskip}{6pt}
\begin{document}
\textbf{2.3\quad The Gaussian integral}

Let $I = \int_{-\infty}^{\infty} e^{-x^2}\,dx$. Squaring $I$ and switching to
polar coordinates $x = r\cos\theta$, $y = r\sin\theta$ gives
\[
I^2 = \int_{0}^{2\pi}\!\int_{0}^{\infty} e^{-r^2}\, r\,dr\,d\theta = \pi ,
\]
so $I = \sqrt{\pi}$. Completing the square in the exponent extends this to
any $a > 0$ and $b \in \mathbb{R}$:
\begin{equation}
\int_{-\infty}^{\infty} e^{-a x^{2}+b x}\,dx=\sqrt{\frac{\pi}{a}}\;e^{\frac{b^{2}}{4 a}}
\tag{2.7}
\end{equation}
Setting $a = \tfrac{1}{2\sigma^2}$ recovers the normalising constant of the
normal density, and differentiating (2.7) with respect to $b$ at $b = 0$
yields every moment $\mathbb{E}[X^{n}]$.
\end{document}
